\begin{frame}<1-3>[t]{新知探索}
\begin{campuscanvas}
% Source slide 6, objects 9; visible 2-3
\only<2-3|handout:1>{\node[anchor=north west,inner sep=0,text=black] at (70.000,150.000) {\begin{minipage}[t]{142.500mm}\raggedright\fontsize{16}{24.00}\selectfont \textbf{2. 集合与元素的关系}\end{minipage}};}
% Source slide 6, objects 8; visible 3
\only<3|handout:1>{\node[anchor=north west,inner sep=0,text=black] at (70.000,260.000) {\begin{minipage}[t]{142.500mm}\raggedright\fontsize{13}{19.50}\selectfont 集合论中最基本的关系是集合和它的元素之间的归属关系。表达归属关系的符号是 $\in$，读作“\red{属于}”。\par 若 $S$ 是一个集合，$a$ 是 $S$ 的一个元素，记作 $a\in S$，读作“$a$ 属于 $S$”。\par 若 $a$ 不是集合 $S$ 的元素，记作 $a\notin S$，读作“$a$ \red{不属于} $S$”。\end{minipage}};}
\end{campuscanvas}
\end{frame}
